\documentclass[11pt,a4paper]{article}

\usepackage[margin=2.4cm]{geometry}
\usepackage{amsmath,amssymb,amsthm}
\usepackage{booktabs}
\usepackage{longtable}
\usepackage{array}
\usepackage{xcolor}
\usepackage[most]{tcolorbox}
\usepackage{fancyhdr}
\usepackage{titlesec}
\usepackage{enumitem}
\usepackage{listings}
\usepackage{hyperref}

\definecolor{strivenavy}{HTML}{1A3D6D}
\definecolor{striveteal}{HTML}{2E8B8B}
\definecolor{strivewarn}{HTML}{C0392B}
\definecolor{strivegrey}{HTML}{F4F6F8}

\hypersetup{colorlinks=true,linkcolor=strivenavy,urlcolor=striveteal,
  pdftitle={STRIVE: Mathematical Specification},
  pdfauthor={Team STRIVE, IISER Bhopal}}

\titleformat{\section}{\Large\bfseries\color{strivenavy}}{\thesection}{0.7em}{}
\titleformat{\subsection}{\large\bfseries\color{striveteal}}{\thesubsection}{0.6em}{}

\pagestyle{fancy}
\fancyhf{}
\lhead{\small\color{strivenavy}STRIVE --- Mathematical Specification}
\rhead{\small\color{strivenavy}SIH26104}
\cfoot{\thepage}
\renewcommand{\headrulewidth}{0.4pt}

\newtcolorbox{whybox}[1]{colback=strivegrey,colframe=striveteal,
  fonttitle=\bfseries,title={Why this choice: #1},boxrule=0.6pt,arc=2pt,left=4pt,right=4pt}
\newtcolorbox{warnbox}[1]{colback=red!3,colframe=strivewarn,
  fonttitle=\bfseries,title={Limitation: #1},boxrule=0.6pt,arc=2pt,left=4pt,right=4pt}

\lstset{basicstyle=\ttfamily\footnotesize,frame=single,framesep=4pt,
  backgroundcolor=\color{strivegrey},rulecolor=\color{gray},breaklines=true}

\newcommand{\code}[1]{\texttt{\small #1}}
\newcommand{\R}{\mathbb{R}}
\newcommand{\Ex}{\mathbb{E}}
\newcommand{\unitvec}[1]{\widehat{#1}}

\begin{document}

\begin{titlepage}
\centering
\vspace*{2cm}
{\Huge\bfseries\color{strivenavy} STRIVE\par}
\vspace{0.6cm}
{\LARGE\color{striveteal} Mathematical Specification\par}
\vspace{0.4cm}
{\large Every formula in the system, and why it was chosen\par}
\vspace{2cm}
{\large\bfseries AI-Powered Real-Time Detection and Prevention\\of Voice Cloning Impersonation Attacks\par}
\vspace{1.5cm}
\begin{tabular}{rl}
\textbf{Problem Statement} & SIH26104\\
\textbf{Organization} & AICTE Cyber Security Cell\\
\textbf{Theme} & Blockchain \& Cybersecurity\\
\textbf{Team} & STRIVE\\
\textbf{Institution} & IISER Bhopal\\
\textbf{Event} & Smart India Hackathon 2026\\
\end{tabular}
\vfill
\begin{tcolorbox}[colback=red!3,colframe=strivewarn,width=0.92\textwidth]
\textbf{Scope of claims.} This document specifies the mathematics that the
implementation actually computes. Where a quantity is a heuristic, a proxy, or
uncalibrated, it is labelled as such. No performance claim appears here that was
not measured; measured values cite the evidence file that contains them.
\end{tcolorbox}
\vspace{1cm}
\end{titlepage}

\tableofcontents
\newpage

%==========================================================
\section{Notation and conventions}

Let $x[n] \in \R$ denote a discrete-time audio signal at sampling rate
$f_s = 16{,}000$\,Hz, normalised so $x[n] \in [-1,1]$. Time in seconds is
$t = n/f_s$.

\begin{longtable}{@{}lp{11.2cm}@{}}
\toprule
\textbf{Symbol} & \textbf{Meaning}\\
\midrule
\endhead
$f_s$ & Sampling rate, fixed at $16{,}000$\,Hz\\
$W$ & Analysis window length in seconds (default $2.0$)\\
$H$ & Hop length in seconds (default $1.0$, supports $0.5$)\\
$N_W = W f_s$ & Window length in samples\\
$N_H = H f_s$ & Hop length in samples\\
$\unitvec{v}$ & Unit-normalised vector, $\unitvec{v} = v/\max(\|v\|_2,\varepsilon)$\\
$s_{\mathrm{glob}}$ & Global retrieval track score (artifact evidence)\\
$s_{\mathrm{sess}}$ & Session consistency track score\\
$s_{\mathrm{coh}}$ & Boundary coherence track score\\
$R_a$ & Authenticity risk (acoustic, temporally smoothed)\\
$R_c$ & Context risk (business metadata)\\
$R_d$ & Decision risk (policy combination)\\
$\varepsilon$ & Numerical floor, $10^{-12}$ unless stated\\
\bottomrule
\end{longtable}

\begin{whybox}{unit normalisation everywhere}
All embeddings are $L^2$-normalised before comparison. On the unit sphere the
inner product equals the cosine similarity,
$\langle \unitvec{u},\unitvec{v}\rangle = \cos\theta$, so a single FAISS inner-product
index computes cosine similarity exactly with no extra work. It also removes
loudness as a confound: a quiet and a loud recording of the same spectral shape map
to the same point.
\end{whybox}

%==========================================================
\section{Signal conditioning}

\subsection{Canonical form and resampling}

Every input is reduced to mono by channel averaging,
$x[n] = \frac{1}{C}\sum_{c=1}^{C} x_c[n]$, then resampled to $f_s$.

Resampling uses a polyphase rational filter. For input rate $r$ and target $f_s$,
let $g=\gcd(r,f_s)$ and set the interpolation and decimation factors
\begin{equation}
L = \frac{f_s}{g}, \qquad M = \frac{r}{g}.
\end{equation}
The signal is upsampled by $L$ (zero insertion), low-pass filtered at
$\omega_c = \pi/\max(L,M)$, and downsampled by $M$.

\begin{whybox}{polyphase over naive interpolation}
Rational $L/M$ resampling with an anti-aliasing filter is exact for the common
telephony rates ($8$\,kHz $\to$ $16$\,kHz is $L{=}2, M{=}1$; $48$\,kHz is
$L{=}1,M{=}3$; $44.1$\,kHz is $L{=}160,M{=}441$). Naive linear interpolation would
alias high-frequency energy into the band where synthesis artifacts live --- precisely
the evidence the detector depends on. Corrupting that band to save compute would be
self-defeating.
\end{whybox}

\subsection{Overlapping window extraction}

Windows are extracted at a fixed hop. Window $k$ spans samples
\begin{equation}
\mathcal{W}_k = \{\,kN_H,\; kN_H+1,\; \ldots,\; kN_H + N_W - 1\,\},
\qquad k = 0,1,2,\ldots
\end{equation}
with start time $t_k = kH$ and end time $t_k + W$. The overlap between consecutive
windows is $W-H$ seconds; the \emph{fresh} (previously unseen) portion is the last
$H$ seconds, except for $k=0$ where the whole window is fresh.

The constraint $0 < H \le W$ is enforced, and both $N_W$ and $N_H$ must be integers.
For $W=2, H=0.5$ the emission rate is $1/H = 2$ windows per second.

\begin{whybox}{fresh-sample accounting}
Voiced-time accumulation uses only the fresh portion:
$T_{\text{voiced}} \mathrel{+}= (|\text{fresh}|/f_s)\cdot a$ where $a$ is the activity
fraction. Using the whole window would count overlapped audio $W/H$ times, inflating
the evidence counter by $4\times$ at a $0.5$\,s hop and letting the system declare
sufficient evidence long before it exists.
\end{whybox}

\subsection{Energy activity gate}

The signal is framed into $20$\,ms blocks ($320$ samples). For block $m$,
\begin{equation}
\rho_m = \sqrt{\frac{1}{320}\sum_{n \in \text{block } m} x[n]^2},
\qquad
a = \frac{1}{M}\sum_{m=1}^{M} \mathbf{1}[\rho_m > \tau_a], \quad \tau_a = 0.008 .
\end{equation}
A window is analysed when $a \ge 0.25$.

\begin{warnbox}{this is an energy detector, not a speech classifier}
$a$ measures whether \emph{anything} is above the noise floor. Measured on 40 matched
ASVspoof clip pairs, additive noise at $10$\,dB SNR raises the mean voiced fraction from
$0.564$ to $1.000$ --- a $1.77\times$ inflation in which every frame is counted as speech.
Consequently noisy clips clear the evidence gate $5.6\times$ more often than clean ones
($28$ vs.\ $5$ clips) while their discriminative power is at its worst. Research mode
substitutes a trained WebRTC VAD. See \S\ref{sec:findings}.
\end{warnbox}

%==========================================================
\section{Spectral representations}

\subsection{Short-time Fourier transform}

With analysis window $w[\cdot]$ of length $N=400$ ($25$\,ms), hop $160$ ($10$\,ms),
and FFT size $512$,
\begin{equation}
X[m,\kappa] = \sum_{n=0}^{N-1} x[mR+n]\, w[n]\, e^{-\mathrm{j}2\pi \kappa n / 512},
\qquad P[m,\kappa] = |X[m,\kappa]|^2 .
\end{equation}
Bin $\kappa$ corresponds to frequency $f_\kappa = \kappa f_s/512 = 31.25\kappa$\,Hz.

\begin{whybox}{25\,ms / 10\,ms framing}
Speech is approximately stationary over $20$--$30$\,ms. Shorter frames lose frequency
resolution ($\Delta f = f_s/512 = 31.25$\,Hz here, enough to resolve the harmonic comb
of a typical $F_0$); longer frames smear transients, which is where vocoder artifacts
concentrate. The $10$\,ms hop gives $2.5\times$ oversampling so a transient cannot fall
between frames.
\end{whybox}

\subsection{Countermeasure embedding (demo extractor)}

The power spectrogram is pooled into $B = 32$ contiguous frequency bands. For band $b$,
\begin{equation}
\beta[m,b] = \frac{1}{|b|}\sum_{\kappa \in b} P[m,\kappa],
\qquad
\tilde\beta[m,b] = \sqrt{\frac{\beta[m,b]}{\sum_{b'} \beta[m,b']}} .
\end{equation}
The square root of the normalised band energy is the \emph{Hellinger} representation:
treating $\beta[m,\cdot]/\sum\beta$ as a probability distribution over bands,
$\|\tilde\beta_1 - \tilde\beta_2\|_2$ is (up to $\sqrt{2}$) the Hellinger distance
between spectra.

The embedding concatenates the temporal mean and standard deviation:
\begin{equation}
v_{\mathrm{cm}} = \unitvec{\big[\; \Ex_m[\tilde\beta[m,\cdot]] \;\big\|\; \sigma_m[\tilde\beta[m,\cdot]] \;\big]} \in \R^{64}.
\end{equation}

\begin{whybox}{mean and variance, not mean alone}
The mean captures the average spectral envelope; the standard deviation captures how
much that envelope \emph{moves}. Synthetic speech is frequently over-smooth in time ---
a vocoder produces less frame-to-frame spectral variance than a physical vocal tract.
The $\sigma$ half of the vector is what lets a purely unsupervised representation
carry any spoof-relevant signal at all.
\end{whybox}

\begin{warnbox}{this is a DSP surrogate}
$v_{\mathrm{cm}}$ is hand-designed, not learned. It is an engineering stand-in so the
streaming architecture can be built and measured before neural weights are available.
Measured AUC on real ASVspoof 2019 LA data is $0.707$ (\S\ref{sec:findings}) --- weak
but non-trivial signal, and explicitly the floor a trained countermeasure must beat.
Every event it produces carries \code{demo\_only: true}.
\end{warnbox}

\subsection{Mel filterbank and cepstral coefficients}

The mel scale approximates human pitch perception:
\begin{equation}
\mathrm{mel}(f) = 2595 \log_{10}\!\left(1 + \frac{f}{700}\right),
\qquad
\mathrm{mel}^{-1}(\mu) = 700\left(10^{\mu/2595} - 1\right).
\end{equation}
Centre frequencies $\{h_i\}_{i=0}^{27}$ are placed uniformly in mel and mapped back to
Hz, giving $26$ triangular filters
\begin{equation}
\Lambda_i(f) = \max\!\left(0,\; \min\!\left(\frac{f-h_i}{h_{i+1}-h_i},\;
\frac{h_{i+2}-f}{h_{i+2}-h_{i+1}}\right)\right).
\end{equation}
MFCCs are the DCT-II of the log filterbank energies, retaining the first $13$:
\begin{equation}
c_j = \sum_{i=0}^{25} \log\!\big(\max(E_i,\varepsilon)\big)\,
\cos\!\left[\frac{\pi j (i + \tfrac12)}{26}\right], \qquad j = 0,\ldots,12 .
\end{equation}

\begin{whybox}{log then DCT}
Source--filter theory models speech as excitation convolved with a vocal-tract
response. Convolution becomes multiplication in the spectral domain and \emph{addition}
after the logarithm, so the log-spectrum separates the two additively. The DCT then
approximately decorrelates the result (it is the KLT for a first-order Markov process)
and compacts energy into low indices, so $13$ coefficients retain most of the
vocal-tract information. Truncating discards the fine harmonic structure, which is
dominated by $F_0$ and by speaker identity rather than synthesis artifacts.
\end{whybox}

\subsection{Spectral shape descriptors}

Per frame, with $\tilde P[\kappa] = P[\kappa] + \varepsilon$:
\begin{align}
\text{centroid:}\quad & \mu_f = \frac{\sum_\kappa f_\kappa \tilde P[\kappa]}{\sum_\kappa \tilde P[\kappa]}
&& \text{spectral ``centre of mass''}\\[2pt]
\text{rolloff:}\quad & f_{85} = \min\Big\{ f_\kappa : \textstyle\sum_{\kappa' \le \kappa} \tilde P[\kappa'] \ge 0.85 \sum_{\kappa'} \tilde P[\kappa'] \Big\}
&& \text{where the energy ends}\\[2pt]
\text{flatness:}\quad & \Phi = \frac{\exp\!\big(\frac{1}{K}\sum_\kappa \ln \tilde P[\kappa]\big)}{\frac{1}{K}\sum_\kappa \tilde P[\kappa]}
&& \text{geometric / arithmetic mean}
\end{align}

\begin{whybox}{flatness as a tonality measure}
By the AM--GM inequality $\Phi \in (0,1]$, with $\Phi = 1$ exactly when the spectrum is
flat (white noise) and $\Phi \to 0$ for a pure tone. It is therefore a scale-invariant
tonality index. Synthesis systems often mismodel the noise/harmonic balance of
fricatives and breathiness, which $\Phi$ summarises in one number.
\end{whybox}

\subsection{Pitch, jitter and shimmer}

Fundamental frequency is estimated by autocorrelation. For a mean-removed frame $f$,
\begin{equation}
r[\ell] = \sum_n f[n]\,f[n+\ell],
\qquad
\ell^\star = \arg\max_{40 \le \ell \le 320} r[\ell],
\qquad
F_0 = \frac{f_s}{\ell^\star}.
\end{equation}
The search range $\ell \in [40,320]$ corresponds to $F_0 \in [50,400]$\,Hz. The
normalised peak $r[\ell^\star]/r[0]$ is the harmonicity; frames below $0.25$ (or $0.5$
for the cycle analysis) are treated as unvoiced and excluded.

With consecutive glottal period lengths $T_1,\ldots,T_K$ and peak-to-peak cycle
amplitudes $A_1,\ldots,A_K$:
\begin{equation}
\text{jitter} = \frac{\frac{1}{K-1}\sum_{i=1}^{K-1} |T_{i+1}-T_i|}{\frac{1}{K}\sum_i T_i},
\qquad
\text{shimmer} = \frac{\frac{1}{K-1}\sum_{i=1}^{K-1} |A_{i+1}-A_i|}{\frac{1}{K}\sum_i A_i}.
\end{equation}

\begin{whybox}{jitter and shimmer as physiological cues}
Human phonation is driven by tissue with mass and elasticity, so consecutive glottal
cycles are never identical: healthy voices exhibit small but non-zero period and
amplitude perturbation. Many synthesis pipelines generate excitation from a smooth
parametric $F_0$ contour and therefore produce jitter and shimmer that are
\emph{too low}. These are dimensionless ratios, so they are invariant to gain and to
the speaker's absolute pitch.
\end{whybox}

\begin{warnbox}{weak, confoundable cues}
Jitter and shimmer are sensitive to background noise, codec quantisation and the
voicing decision itself. They are supporting evidence only and are never used as a
standalone spoofing verdict.
\end{warnbox}

%==========================================================
\section{Channel intelligence}
\label{sec:channel}

Channel measurements describe the \emph{transport}, never the speaker. A degraded
channel means our evidence deserves less trust; it is never itself evidence of a clone.
Every quantity below is deterministic and model-free. Unavailable quantities are
returned as \code{null}, never as $0$.

\subsection{Clipping ratio and level}
\begin{equation}
\mathrm{clip} = \frac{1}{N}\sum_{n} \mathbf{1}\big[\,|x[n]| \ge 0.99\,\big],
\qquad
L_{\mathrm{dBFS}} = 20\log_{10}\!\sqrt{\tfrac{1}{N}\textstyle\sum_n x[n]^2}.
\end{equation}
$L_{\mathrm{dBFS}}$ is \code{null} when the RMS is below $10^{-9}$ (digital silence),
since $\log_{10} 0$ is undefined and reporting $0$\,dBFS would falsely indicate a
full-scale signal.

\subsection{Occupied bandwidth}

Let $\bar P[\kappa]$ be the time-averaged power spectrum. The occupied bandwidth is the
$99\%$ cumulative-power rolloff:
\begin{equation}
B = f_{\kappa^\star}, \qquad
\kappa^\star = \min\Big\{\kappa : \textstyle\sum_{\kappa' \le \kappa}\bar P[\kappa'] \ge 0.99 \sum_{\kappa'}\bar P[\kappa'] \Big\}.
\end{equation}

\begin{whybox}{cumulative rolloff over a $-3$\,dB cutoff}
A $-3$\,dB cutoff must be measured relative to a passband reference, which is unstable
for speech because the passband level varies with phonation. The cumulative-energy
definition needs no reference level and degrades gracefully: it always returns a
frequency, and the fraction ($0.99$) sets exactly how much tail energy is ignored.
Measured on the codec matrix it cleanly separates narrowband transport
($\approx 2672$\,Hz for G.711 and $8$\,kHz resampling) from clean wideband
($\approx 3625$\,Hz).
\end{whybox}

\begin{warnbox}{content dependence}
$B$ depends on what is being said, not only on the channel. A low-pitched or quiet
speaker on a wideband line can read as narrowband. $B$ is a channel \emph{hint}, not a
codec classifier.
\end{warnbox}

\subsection{Signal-to-noise ratio proxy}

Noise is estimated per frequency bin by \emph{minimum statistics}: over the frame index
$m$, the low quantile of the power in each bin approximates the stationary noise floor,
because speech is intermittent while noise is persistent.
\begin{equation}
\hat n[\kappa] = Q_{0.10}\big(\{P[m,\kappa]\}_m\big),
\qquad
\hat N = \sum_\kappa \hat n[\kappa],
\qquad
\hat S = \sum_\kappa \Ex_m\big[P[m,\kappa]\big] - \hat N,
\end{equation}
\begin{equation}
\mathrm{SNR}_{\mathrm{dB}} = \mathrm{clip}\!\left(10\log_{10}\frac{\hat S}{\max(\hat N,\; 10^{-6}\hat T)},\; 0,\; 60\right),
\qquad \hat T = \hat S + \hat N .
\end{equation}

\begin{whybox}{per-bin minimum statistics, not frame energy percentiles}
The first implementation used the ratio of the $95$th to the $10$th percentile of
\emph{frame energy}. That collapses on stationary input: for a signal with constant
energy every frame is identical, the ratio is $1$, and the reported SNR is $0$\,dB
regardless of how clean the signal is. In testing, a clean tone scored $0.0$\,dB while
the same tone plus noise scored $0.6$\,dB --- the metric was inverted. Operating per
frequency bin fixes this, because even a stationary tone occupies few bins and leaves
the rest at the true noise floor. The floor $10^{-6}\hat T$ prevents a division blow-up
on digitally clean input, which saturates to the $60$\,dB cap instead of returning
``unmeasurable''.
\end{whybox}

\begin{warnbox}{a proxy, not a measurement}
Verified monotonic in added noise (true $0/5/10/20/40$\,dB $\to$ estimated
$11.8/15.1/19.5/28.7/47.9$\,dB) but biased high by roughly $8$\,dB, because the $10$th
percentile underestimates the floor. Use it for ordering and for reliability weighting;
do not quote it as a calibrated SNR.
\end{warnbox}

\subsection{Cross-window discontinuity}

Let $\phi_k = (\rho_k, \mu_{f,k}, \nu_k, e_k)$ summarise fresh audio in window $k$:
RMS, spectral centroid, noise floor ($10$th percentile frame energy) and final sample.
Define a scale-free log jump
\begin{equation}
J(u,v;\delta) = \mathrm{clip}\!\left(\frac{\big|\ln\frac{u+\delta}{v+\delta}\big|}{\ln 10},\,0,\,1\right),
\end{equation}
which reaches $1$ at one order of magnitude of change. The discontinuity score is
\begin{equation}
D_k = \max\Big\{
J(\rho_k,\rho_{k-1};10^{-3}),\;
J(\nu_k,\nu_{k-1};10^{-8}),\;
\tfrac{|\mu_{f,k}-\mu_{f,k-1}|}{f_s/4},\;
\tfrac{|x_k[0]-e_{k-1}|}{\max(4\rho_k,\,10^{-6})}
\Big\}
\end{equation}
with $D_0 = 0$ (no predecessor).

\begin{whybox}{maximum rather than sum}
A splice need only disturb \emph{one} of these quantities: an edit can preserve energy
while shifting the noise floor, or preserve spectrum while leaving a sample-level step.
A sum would dilute a single strong indicator across four terms; the maximum preserves
it. The additive offsets $\delta$ keep the log finite when either side is silent.
\end{whybox}

\begin{warnbox}{packet loss is invisible here}
Measured on the codec matrix, $1\%$, $3\%$ and $5\%$ packet loss produce channel
readings identical to clean audio. Twenty-millisecond zero-erasures do not move the
per-window aggregate energy, centroid or noise floor enough to register. Consequently
\code{packet\_loss\_rate} must be supplied by the transport and is never inferred from
the waveform. Detecting erasures from audio requires a dedicated zero-run detector, not
a refinement of $D_k$.
\end{warnbox}

\subsection{Reliability aggregation}

Channel quality is one minus a sum of bounded penalties:
\begin{equation}
Q = \mathrm{clip}\!\left(1 - \sum_{i} w_i \,\pi_i,\; 0,\; 1\right),
\end{equation}
\begin{align}
\pi_{\mathrm{clip}} &= \min\!\left(1, \tfrac{\mathrm{clip}}{0.02}\right), &
\pi_{\mathrm{disc}} &= D_k, \\
\pi_{\mathrm{snr}} &= \mathrm{clip}\!\left(\tfrac{25 - \mathrm{SNR_{dB}}}{25-5},0,1\right), &
\pi_{\mathrm{bw}} &= \mathrm{clip}\!\left(\tfrac{7000 - B}{7000-3400},0,1\right),
\end{align}
with default weights $w = (0.25, 0.20, 0.30, 0.25)$. A penalty whose input is
unavailable is \emph{omitted}, not set to zero, and $Q$ itself is \code{null} when the
window contains no usable audio.

\begin{whybox}{linear ramps with explicit anchors}
Each $\pi_i$ is a clipped linear ramp between a ``good'' and a ``bad'' anchor
($25 \to 5$\,dB SNR; $7000 \to 3400$\,Hz bandwidth, the latter being the G.711
passband). Linear ramps are monotone, bounded, and every constant is a physically
interpretable number that a reviewer can argue with --- unlike a fitted sigmoid, which
would imply calibration we have not earned. All constants live in configuration.
\end{whybox}

\begin{tcolorbox}[colback=strivegrey,colframe=strivenavy,title=\textbf{$Q$ is reliability, not risk}]
$Q$ answers ``how much should we trust the acoustic evidence from this window?'' It
never answers ``is this a clone?''. A pristine channel does not make a caller genuine,
and a bad channel does not make a caller fake. Keeping these separate is why $Q$ is
consumed by fusion weighting and never added to the risk score.
\end{tcolorbox}

%==========================================================
\section{Evidence tracks}

\subsection{Track 1 --- labelled-reference retrieval}

The reference index holds unit embeddings $\{\unitvec{u}_i\}_{i=1}^{M}$ with labels
$y_i \in \{0,1\}$ ($0$ bonafide, $1$ spoof). For a query $\unitvec{q}$, retrieve the $k$
nearest neighbours by inner product,
$\mathcal{N}_k(\unitvec{q}) = \operatorname{arg\,top-}k_i \langle \unitvec{q},\unitvec{u}_i\rangle$, and score
\begin{equation}
s_{\mathrm{glob}} = \frac{1}{k}\sum_{i \in \mathcal{N}_k(\unitvec{q})} y_i \;\in [0,1],
\end{equation}
the spoof fraction of the neighbourhood. The index is partitioned by language, falling
back to a global pool when a partition has fewer than $20$ entries or contains a single
class; the fallback is reported in the event as a reason code.

\begin{whybox}{$k$-NN retrieval over a trained classifier head}
A retrieval score is interpretable and auditable: every decision can be traced to the
specific labelled clips that produced it, which matters for a fraud-prevention tool
that must justify holding a transaction. It also extends to new attack generators by
adding data rather than retraining. The cost is that quality is bounded by the corpus:
a single-class or tiny partition cannot produce a meaningful score, which is exactly
why the guard returns \code{null} instead of a number.
\end{whybox}

\subsection{Track 2 --- session consistency}

Within a call, a bounded store of accepted profile vectors is maintained. The similarity
of the current window to that store is
\begin{equation}
\sigma = \tfrac12 \underbrace{\Big(\tfrac{1}{k}\!\!\sum_{i \in \mathcal{N}_k}\!\! \langle \unitvec{p},\unitvec{p}_i\rangle\Big)}_{\text{speaker-level profile}}
\;+\; \tfrac12 \underbrace{\Big(\tfrac{1}{|S|}\sum_{s \in S} \tfrac{1}{k}\!\!\sum_{j \in \mathcal{N}_k(s)}\!\! \langle \unitvec{s},\unitvec{s}_j\rangle\Big)}_{\text{segment-level}},
\qquad
s_{\mathrm{sess}} = \mathrm{clip}(1-\sigma,\,0,\,1).
\end{equation}

\paragraph{Trust-gated admission.} A window updates the store only if
\begin{equation}
s_{\mathrm{glob}} < 0.30 \quad\text{and}\quad \sigma > 0.70 .
\end{equation}
Bootstrap requires \emph{all} initial window scores below $0.30$, at least $3$ valid
windows, and $\ge 4$\,s of voiced audio by the $5$\,s mark; otherwise the store is
permanently blocked for that call.

\begin{whybox}{a dual gate against profile poisoning}
The session profile answers ``is the voice still the same as it was?'' --- which is only
meaningful if the initial voice was genuine. Without gating, an attacker who is cloned
from the first second would have the clone admitted as the reference, and every
subsequent cloned window would score as perfectly consistent. Requiring low global risk
\emph{and} high self-similarity means an admitted window must look genuine to the
independent retrieval track and be consistent with what is already stored.
\end{whybox}

\begin{warnbox}{not proof against slow poisoning}
The gate stops obvious contamination. A patient attacker who drifts the voice slowly,
keeping each step inside both thresholds, is not addressed by this mechanism and has not
been evaluated.
\end{warnbox}

\subsection{Track 3 --- boundary coherence}

The seam between old and new audio lies at $\tau = W - H$ seconds inside the current
window. Among adjacent segment pairs $(a,b)$ bracketing $\tau$ with a gap
$b_{\text{start}} - a_{\text{end}} \le 0.12$\,s, select the pair whose midpoint is
closest to $\tau$, and require that midpoint to be within $0.2$\,s of $\tau$. Then
\begin{equation}
s_{\mathrm{coh}} = \mathrm{clip}\big(1 - \langle \unitvec{v}_a, \unitvec{v}_b \rangle,\; 0,\; 1\big).
\end{equation}
If no admissible pair exists, the track returns \code{null}.

\begin{whybox}{clipping a quantity whose natural range is $[0,2]$}
Cosine similarity lies in $[-1,1]$, so $1-\cos$ lies in $[0,2]$. Values above $1$ mean
anti-correlated embeddings, which for spectral segments indicates a pathological
comparison rather than a stronger splice. Clipping keeps all track scores on a common
$[0,1]$ scale so the fusion weights mean the same thing for every track.
\end{whybox}

\begin{warnbox}{a weak and confoundable cue}
Adjacent \emph{natural} phonemes legitimately have different embeddings --- that is what
makes them different phonemes. High $s_{\mathrm{coh}}$ therefore does not imply
splicing, and the cue is further confounded by packet loss and phonetic content. It is
supporting evidence with a small weight and has not been validated as a discriminator.
\end{warnbox}

%==========================================================
\section{Fusion and temporal smoothing}

\subsection{Age-scheduled weights}

Let $A$ be the call age in seconds. The nominal weight vector over
$(s_{\mathrm{glob}}, s_{\mathrm{sess}}, s_{\mathrm{coh}})$ is
\begin{equation}
w(A) = \begin{cases}
(0.80,\; 0.00,\; 0.20) & A < 15\\
(0.50,\; 0.30,\; 0.20) & 15 \le A \le 60\\
(0.30,\; 0.50,\; 0.20) & A > 60 .
\end{cases}
\end{equation}

\begin{whybox}{shifting weight from global to session evidence}
Early in a call there is no session history, so the session weight is exactly zero ---
not a small number, zero. As the call proceeds an in-call profile accumulates and
becomes informative about mid-call speaker substitution, which the global track cannot
see. The schedule encodes that changing information content.
\end{whybox}

\subsection{Availability masking}

Let $m_i = \mathbf{1}[s_i \ne \code{null}]$. The effective weights are
\begin{equation}
\tilde w_i = \frac{m_i\, w_i}{\sum_j m_j w_j}
\quad\text{if}\quad \textstyle\sum_j m_j w_j > 0,
\end{equation}
and the raw fused score is $s_{\mathrm{raw}} = \sum_i \tilde w_i s_i$. If
$s_{\mathrm{glob}}$ is unavailable, \emph{no} risk is published.

\begin{tcolorbox}[colback=strivegrey,colframe=strivenavy,title=\textbf{Absent evidence is not evidence of innocence}]
The single most important line in the fusion is that a missing track is masked and the
remaining weights renormalised --- it does \emph{not} contribute $s_i = 0$. Substituting
zero would let a failed model, a cold-start session profile, or a disabled identity
branch actively push the risk \emph{down}, manufacturing confidence from silence. The
same principle makes a stale scheduled branch unavailable rather than reusing an expired
value indefinitely.
\end{tcolorbox}

\subsection{Exponential moving average}

\begin{equation}
R_a^{(k)} = \alpha R_a^{(k-1)} + (1-\alpha)\, s_{\mathrm{raw}}^{(k)},
\qquad \alpha = 0.7, \qquad R_a^{(-1)} = 0 .
\end{equation}
This is a first-order IIR low-pass filter with impulse response
$(1-\alpha)\alpha^j$, effective memory
$\sum_j j(1-\alpha)\alpha^j = \alpha/(1-\alpha) \approx 2.33$ windows, and
$-3$\,dB cutoff at $\omega_c = \arccos\!\big(2 - \tfrac{1+\alpha^2}{2\alpha}\big)$.

Starting from zero, $k$ consecutive maximal inputs give
\begin{equation}
R_a^{(k)} = 1 - \alpha^{k+1}.
\end{equation}

\begin{whybox}{smoothing is a false-positive control}
A single anomalous window --- a cough, a codec glitch, one bad retrieval --- must not
trigger a transaction hold. Requiring sustained evidence is the mechanism that makes
``verify, do not block'' workable in production.
\end{whybox}

\begin{warnbox}{documented counterexamples, deliberately preserved}
Two consequences are pinned by unit tests rather than tuned away:
\begin{enumerate}[nosep,leftmargin=1.4em]
\item \textbf{Warm-up.} Three maximal updates give $1-0.7^3 = 0.657 < 0.75$, so an alert
cannot fire in three windows; five give $1-0.7^5 \approx 0.832$. The source design
document's ``alert in three chunks'' claim is arithmetically wrong.
\item \textbf{Mature-call dilution.} For $A > 60$ with all tracks available,
$s = (1,0,0)$ fuses to $0.30$ --- below the warning band. A static clone that matches
its own bootstrap profile can therefore suppress a maximal global alarm. This is a real
weakness of the specified schedule and is recorded, not hidden.
\end{enumerate}
\end{warnbox}

%==========================================================
\section{Risk separation and policy}

\subsection{Context risk}

Business metadata maps to a bounded score by additive contributions:
\begin{equation}
R_c = \min\Big(1,\;
\underbrace{\begin{cases}0.35 & \text{amount} \ge 10^6\\ 0.15 & \text{amount} \ge 10^5\\ 0 & \text{otherwise}\end{cases}}_{\text{transaction value}}
+\; 0.20\,\mathbf{1}[\text{urgent}]
+\; 0.25\,\mathbf{1}[\text{new beneficiary}]
+\; 0.20\,\mathbf{1}[\text{privileged}]\Big).
\end{equation}

\subsection{Decision risk}

\begin{equation}
\boxed{\;R_d = 1 - (1-R_a)\left(1 - \tfrac12 R_c\right)\;}
\end{equation}

\begin{whybox}{a noisy-OR, and why context is halved}
Reading $R_a$ and $R_c$ as independent probabilities that \emph{something is wrong},
$1-R_d = (1-R_a)(1-\tfrac12 R_c)$ is the probability that neither concern fires. This is
the noisy-OR combination: it is monotone in both arguments, saturates correctly at $1$,
and never lets a low context risk \emph{reduce} the acoustic risk --- since
$R_d \ge R_a$ always, with equality when $R_c = 0$. The factor $\tfrac12$ caps how much
business metadata alone can contribute, because a large payment is not evidence of voice
cloning. A wire transfer of INR $10^7$ with pristine audio reaches only
$R_d = 0.175$.
\end{whybox}

\begin{tcolorbox}[colback=strivegrey,colframe=strivenavy,title=\textbf{The separation invariant}]
$R_a$ is a function of acoustic evidence \emph{only}. Context can never alter it. The
two combine at exactly one place, $R_d$, and a unit test asserts that changing context
metadata leaves $R_a$ bit-identical. Collapsing the three numbers into one would make it
impossible for an operator to know whether a call was flagged because the voice was
suspicious or merely because the payment was large.
\end{tcolorbox}

\subsection{Decision bands}

With warning threshold $\theta_w = 0.50$ and alert threshold $\theta_a = 0.75$:
\begin{equation}
\text{state} =
\begin{cases}
\textsc{analyzing} & R_a = \code{null}\\
\textsc{critical/high} & R_d \ge \theta_a\\
\textsc{review} & R_d \ge \theta_w \;\text{ or }\; R_c \ge 0.7\\
\textsc{low} & \text{otherwise.}
\end{cases}
\end{equation}
A publishable risk additionally requires $T_{\text{voiced}} \ge 4$\,s; otherwise the
state is \textsc{analyzing} with reason \code{LOW\_EVIDENCE}. Alerts \emph{latch} until
an independent verification is recorded.

\begin{whybox}{latching, and never publishing a green state early}
Risk is non-monotone: the age schedule can lower a score after a genuine alarm, so an
unlatched policy would silently release a hold. Latching makes the safe action sticky
until a human verification clears it. Similarly, the evidence gate ensures silence or a
cold start reports \textsc{analyzing} rather than \textsc{low} --- an unknown is never
displayed as a clean bill of health.
\end{whybox}

%==========================================================
\section{Scheduling and latency}

\subsection{Multi-rate branch cadence}

Branch $i$ with cadence $C_i$ runs at window $k$ (audio time $t_k$) iff
$t_k - t_{\text{last},i} \ge C_i$. A published result expires after $3C_i$; between runs
the cached result is reused. Results are in one of three states --- missing, fresh, or
stale --- and the latter two are never conflated with a zero score.

\begin{whybox}{cadence on the audio clock, not the wall clock}
If cadence were measured in wall-clock time, an accelerated offline replay (which
processes $40$\,s of audio in $0.25$\,s) would find every branch ``not yet due'' and skip
nearly all of them, making benchmarks meaningless and irreproducible. Measuring on the
audio timeline makes an offline replay schedule branches exactly as a live call would,
so evaluation and production agree by construction.
\end{whybox}

Measured effect at $W=2$, $H=0.5$: the artifact branch runs $39$ times instead of $77$
while all $77$ windows still emit a fused score, giving a $1.63\times$--$1.67\times$
throughput gain reproduced on two different processors.

\subsection{Capture backpressure}

The capture queue holds at most $N_q = 8$ completed windows. On overflow the
\emph{oldest} window is discarded and counted.

\begin{whybox}{drop oldest, not newest}
A live risk score is only useful if it describes the audio the caller is speaking now.
Discarding the newest window would freeze the operator's view in the past while the
backlog drains --- the precise failure the score exists to prevent. Dropping the oldest
bounds latency at $N_q H$ seconds and makes the loss explicit rather than letting delay
grow without limit. A drop also resets the continuity state, so the gap is not
mis-scored as a splice.
\end{whybox}

\subsection{Percentiles and real-time factor}

The nearest-rank percentile of a sorted sample $z_{(1)} \le \cdots \le z_{(n)}$ is
\begin{equation}
P_p = z_{(\lceil pn/100 \rceil)},
\qquad
\mathrm{RTF} = \frac{\text{audio seconds processed}}{\text{compute seconds}} .
\end{equation}

\begin{whybox}{nearest-rank, and why $p95$ rather than the mean}
Nearest-rank returns an \emph{observed} value, never an interpolation between two
measurements that were never made. Latency distributions are right-skewed, so the mean
hides the tail that actually breaks a real-time budget; $p95$ against the hop $H$ is the
meaningful pass/fail test. $\mathrm{RTF} > 1$ is the condition for keeping up with live
audio. An empty sample returns \code{null}, not $0$, since $0$\,ms would read as
``instantaneous''.
\end{whybox}

%==========================================================
\section{Evaluation metrics}

Let scores $\{z_i\}$ and labels $\{y_i\}$, $y_i = 1$ for spoof. For threshold $\theta$,
\begin{equation}
\mathrm{FPR}(\theta) = \frac{|\{i: y_i=0,\, z_i \ge \theta\}|}{|\{i : y_i = 0\}|},
\qquad
\mathrm{TPR}(\theta) = \frac{|\{i: y_i=1,\, z_i \ge \theta\}|}{|\{i : y_i = 1\}|}.
\end{equation}

\paragraph{Equal error rate.} With $\mathrm{FNR} = 1-\mathrm{TPR}$, the EER is the value
where $\mathrm{FPR} = \mathrm{FNR}$. Since the empirical ROC is a step function the
crossing generally falls between two operating points $a,b$; the implementation
interpolates linearly on $g(\theta) = \mathrm{FPR}-\mathrm{FNR}$:
\begin{equation}
\lambda = \frac{-g_a}{g_b - g_a},
\qquad
\mathrm{EER} = \mathrm{FPR}_a + \lambda\,(\mathrm{FPR}_b - \mathrm{FPR}_a).
\end{equation}

\paragraph{ROC-AUC.} $\mathrm{AUC} = \Pr\big(z_i > z_j \mid y_i = 1, y_j = 0\big)$,
the probability that a random spoof outranks a random bonafide --- equivalently the
normalised Mann--Whitney $U$ statistic. It is threshold-free.

\paragraph{Brier score and calibration error.}
\begin{equation}
\mathrm{BS} = \frac{1}{n}\sum_i (z_i - y_i)^2,
\qquad
\mathrm{ECE} = \sum_{b=1}^{10} \frac{|B_b|}{n}\,\big|\,\bar z_{B_b} - \bar y_{B_b}\,\big|.
\end{equation}

\begin{whybox}{reporting a family of metrics, not a single headline}
EER assumes the two error types are equally costly, which is false here: a false
positive costs a callback, a false negative costs a fraudulent transfer. AUC is
threshold-free but insensitive to calibration. Brier and ECE ask whether a score of
$0.8$ actually means $80\%$, which matters because the policy engine thresholds on the
number. Together they describe ranking quality, operating-point behaviour, and
calibration --- no one of them is sufficient.
\end{whybox}

\begin{warnbox}{low-FPR operating points need cohort size}
With $n_0$ genuine clips the finest resolvable false-positive rate is $1/n_0$. The
measured cohort has $n_0 = 65$, so $1/65 \approx 1.5\%$ and a ``TPR at $1\%$ FPR''
figure is \emph{not resolvable}. The tooling emits \code{low\_fpr\_sample\_warning} when
$n_0 < 100$ and that number must not be quoted.
\end{warnbox}

%==========================================================
\section{Measured results and findings}
\label{sec:findings}

Corpus: ASVspoof 2019 LA evaluation partition (ODC-BY), $150$ bonafide / $150$ spoof,
speaker-disjoint $31/31$ split, reference index built from clean audio only. Extractor:
\code{dsp-surrogate-v1}. Full report in \code{evidence/codec-benchmark.json}.

\begin{center}
\begin{tabular}{@{}lrrrrrr@{}}
\toprule
\textbf{Condition} & \textbf{EER \%} & \textbf{$\Delta$EER} & \textbf{AUC} & \textbf{$\Delta$AUC} & \textbf{Bona.\ med.} & \textbf{Spoof med.}\\
\midrule
\texttt{clean} & 37.01 & --- & 0.707 & --- & 0.350 & 0.463\\
\texttt{opus\_16k} & 37.43 & $+0.42$ & 0.688 & $-0.019$ & 0.400 & 0.475\\
\texttt{opus\_32k} & 39.66 & $+2.66$ & 0.686 & $-0.021$ & 0.400 & 0.500\\
\texttt{narrowband\_8k} & 37.40 & $+0.40$ & 0.676 & $-0.031$ & 0.275 & 0.325\\
\texttt{g711\_ulaw} & 37.49 & $+0.48$ & 0.676 & $-0.031$ & 0.275 & 0.325\\
\textbf{\texttt{noise\_10db}} & \textbf{42.14} & $\mathbf{+5.13}$ & \textbf{0.639} & $\mathbf{-0.068}$ & 0.450 & 0.450\\
\bottomrule
\end{tabular}
\end{center}

\paragraph{Finding 1 --- noise dominates codec compression.}
Every codec leg costs only $0.019$--$0.031$ AUC. Additive noise at $10$\,dB SNR costs
$0.068$, more than twice the worst codec, and collapses the median separation between
bonafide and spoof to \emph{exactly zero}. For this feature set the dominant robustness
threat is noise, not compression.

\paragraph{Finding 2 --- noise inflates perceived evidence while destroying it.}
Clips clearing the live-call evidence gate: $5$ (clean) versus $28$ (\code{noise\_10db}),
a $5.6\times$ increase. The mechanism is the energy activity gate of \S2.3, whose voiced
fraction rises from $0.564$ to $1.000$ under the same condition. The system becomes more
confident it has sufficient evidence exactly when that evidence is least discriminative.
Since $\mathrm{SNR_{dB}}$ already reads $19.6$\,dB and $Q$ is available per window, the
fusion has the signal required to discount this evidence; wiring $Q$ into the fusion
weights is the outstanding work.

\begin{tcolorbox}[colback=red!3,colframe=strivewarn,title=\textbf{Interpretation}]
$\mathrm{AUC} = 0.707$ is weak but non-trivial signal from an unsupervised DSP
representation on real spoofed speech. It is \emph{not} a usable detector and not a
deepfake verdict. Its purpose is to establish a measured floor on a real corpus, so that
a trained countermeasure can be compared against something rather than asserted. No
value in this document was estimated, extrapolated, or tuned to look better.
\end{tcolorbox}

%==========================================================
\section{Summary of design principles}

\begin{enumerate}[leftmargin=1.6em]
\item \textbf{Absent evidence is masked, never zeroed.} Missing, stale and failed
branches are excluded from a renormalised weighted sum instead of contributing a score
that would read as ``genuine''.
\item \textbf{Unavailable measurements are \code{null}.} A quantity that could not be
computed is never reported as $0$, because $0$ is a legitimate value for most of these
quantities and would be indistinguishable from a real measurement.
\item \textbf{Channel reliability is orthogonal to risk.} $Q$ scales how much the
acoustic evidence is trusted; it is never added to the risk.
\item \textbf{Authenticity and context never mix except at the decision.} Enforced by a
unit test.
\item \textbf{Bounded ranges throughout.} Every track score, penalty and risk is clipped
to $[0,1]$ so weights carry consistent meaning.
\item \textbf{Scale-free comparisons.} Log ratios, cosine similarity, and dimensionless
ratios (jitter, shimmer, flatness) so gain and pitch do not act as confounds.
\item \textbf{Counterexamples are pinned, not tuned.} The EMA warm-up and mature-call
dilution weaknesses are asserted by tests so they cannot be silently ``fixed'' by moving
a threshold.
\end{enumerate}

\vspace{0.8cm}
\begin{center}
\color{strivenavy}\itshape
Team STRIVE --- IISER Bhopal --- Smart India Hackathon 2026\\
\normalfont\small Secure Voices $\cdot$ Stronger India
\end{center}

\end{document}
